\par\noindent\SourceLeaf{54}%
Setzen wir in diesem Ausdruck $(2q-1) = r$, so ergibt sich durch diese Substitution der Koeffizient von $\cos r\varphi$ (bei ungeradem $r$) in einer Form mit dem Anfangsglied

\[
= \frac{1}{2^{r-1}} {}^{\nu}\mathfrak{B}^{(r)} n^{r},
\]

\noindent und dem $h^{\text{tus}}$ Glied

\[
= \frac{1}{2^{r+2h-1}} {}^{r+2h}\mathfrak{B}^{(r+h)}  .  {}^{\nu}\mathfrak{B}^{(r+2h)} n^{r+2h}.
\]

\noindent Ferner war das Anfangsglied des zu $\cos 2q\varphi$ gehörenden Koeffizienten

\[
\frac{1}{2^{2q-1}} {}^{\nu}\mathfrak{B}^{(2q)} n^{2q}, \qquad h^{\text{tus}} = \frac{1}{2^{2q+2h-1}} {}^{2q+2h}\mathfrak{B}^{(2q+h)}  .  {}^{\nu}\mathfrak{B}^{(2q+2h)} n^{2q+2h}.
\]

\noindent Setzen wir auf dieselbe Weise $2q = r$, so erhalten wir durch diese Substitution für den Koeffizienten von $\cos r\varphi$ (bei geradem $r$) einen Ausdruck mit dem Anfangsglied

\[
= \frac{1}{2^{r-1}} {}^{\nu}\mathfrak{B}^{(r)} n^{r},
\]

\noindent und dem allgemeinen Glied $h^{\text{tus}}$

\[
= \frac{1}{2^{r+2h-1}} {}^{r+2h}\mathfrak{B}^{(r+h)}  .  {}^{\nu}\mathfrak{B}^{(r+2h)} n^{r+2h}.
\]

\noindent Ob wir $r$ gerade oder ungerade wählen, die Ausdrücke, die wir für den Koeffizienten von $\cos r\varphi$ erhalten, sind identisch. Das zeigt, dass für alle Glieder von $(1+n\cos\varphi)^{\nu}$ ein und dasselbe Gesetz gilt, das wir daher in der folgenden Formel knapper angeben.

Unsere Größe $(1+n\cos\varphi)^{\nu}$ ist also gleich der Reihe

\[
A + A^{(1)}\cos\varphi + A^{(2)}\cos 2\varphi + .\,.\,.\,. + A^{(r)}\cos r\varphi + .\,.\,.
\]

\noindent mit folgendem Gesetz für die Koeffizienten. Der Anfangskoeffizient $A$ ist gleich einer Reihe mit dem Anfangsglied $=1$ und dem allgemeinen Glied $h^{\text{tus}}$

\[
= \frac{1}{2^{2h}} {}^{2h}\mathfrak{B}^{(h)}  .  {}^{\nu}\mathfrak{B}^{(2h)} n^{2h};
\]

\noindent Der beliebige $r^{\text{tus}}$ Koeffizient $A^{(r)}$ ist durch eine Reihe mit dem Anfangsglied

\[
= \frac{1}{2^{r-1}} {}^{\nu}\mathfrak{B}^{(r)} n^{r};
\]

\noindent und dem folgenden Koeffizienten des $h^{\text{ti}}$ Gliedes derselben Reihe gegeben

\[
= \frac{1}{2^{r+2h-1}} {}^{\nu}\mathfrak{B}^{(r+2h)}  .  {}^{r+2h}\mathfrak{B}^{(r+h)} n^{r+2h}.
\]
